\documentclass{article}

% NeurIPS 2025 style. Use \usepackage[preprint]{neurips_2025} for a preprint,
% or \usepackage[final]{neurips_2025} for a camera-ready version.
\usepackage[preprint]{neurips_2025}

\usepackage[utf8]{inputenc}
\usepackage[T1]{fontenc}
\usepackage{hyperref}
\usepackage{url}
\usepackage{booktabs}
\usepackage{amsfonts}
\usepackage{nicefrac}
\usepackage{microtype}
\usepackage{xcolor}
\usepackage{graphicx}
\usepackage{siunitx}

\title{When the Ceiling Is the Data: A Falsification-First Audit of\\
Wearable Stress Detection in Nurses}

\author{%
  Elvis Han\thanks{Equal contribution.} \\
  Johns Hopkins University \\
  \texttt{ehan20@jh.edu} \\
  \And
  Eric Guan\footnotemark[1] \\
  Johns Hopkins University \\
  \texttt{guan8zhi@gmail.com} \\
}

\begin{document}

\maketitle

\begin{abstract}
Wearable stress detection is usually reported as a classification score. We
argue that on small, self-reported, single-site datasets the more useful
question is whether the ceiling lies in the method or in the data, and we
answer it by elimination. Using a public dataset of 15 nurses monitored for one
week each during the COVID-19 pandemic, we first establish that the labels mark
a physiologically distinguishable state: skin conductance rises at reported
onset in 64.3\% of episodes (sign test $p = 0.0009$), while heart rate is
indistinguishable from a coin flip (48.6\%, $p = 0.80$). We then construct a
negative class, which the dataset does not contain, by activity-matched
sampling from unlabelled time, and show through falsification probes that the
resulting classifier is neither an activity detector (AUC 0.51) nor a
shift-schedule detector (AUC 0.51). A gradient-boosting baseline reaches
AUC 0.765 but recovers only 10.6\% of episodes at one false alarm per worn
hour, and two of ten subjects are undetectable at any usable threshold. Three
distinct treatments of the labels, including non-negative positive-unlabelled
risk estimation across a bracketed class prior, all fail to improve on this.
We report two findings that run against expectation: careful negative
construction improves rather than inflates performance, and a widely quoted
accuracy figure on this dataset is difficult to reconcile with
subject-grouped evaluation. We conclude that the binding constraint is the
label supply, not the model, and we quantify what a future collection would
need.
\end{abstract}

\section{Introduction}

Occupational stress among nurses is associated with burnout, turnover, and
adverse patient events. Wrist-worn sensors offer a continuous and objective
alternative to retrospective self-report, and a public dataset now exists for
studying it: Hosseini et al.\ \cite{hosseini2022} monitored 15 nurses for one
week each during the COVID-19 outbreak, recording electrodermal activity,
heart rate, skin temperature, and acceleration, with periodic
smartphone-administered stress reports.

Published results on this dataset are strong. Mathur et al.\
\cite{mathur2024} report a weighted $F_1$ of 0.99 for two-level stress
classification. We were unable to reproduce anything close to this, and our
audit suggests why: the dataset contains 178 high-stress episodes from
15 people, adjacent analysis windows are near-duplicates of one another, and
no negative class was ever recorded.

This paper takes a different posture. Rather than reporting a score, we treat
the study as a sequence of eliminations against one question:

\begin{quote}
\emph{Is the ceiling on this task set by the method or by the data?}
\end{quote}

Each phase of the analysis removes one candidate explanation. Section
\ref{sec:labels} asks whether the labels mark anything physiological at all.
Section \ref{sec:negatives} asks whether the constructed comparison group leaks
a confound. Section \ref{sec:baseline} asks how far standard features and a
standard learner get. Section \ref{sec:pu} asks whether a different treatment
of the labels helps. What survives is an argument about the data.

\paragraph{Contributions.}
\begin{enumerate}
\item An event-triggered analysis establishing that only one of four channels
      responds at reported onset, which reorders the feature priorities for
      this dataset and, we suspect, for wrist-worn stress detection generally.
\item A negative-class construction with falsification probes, and the finding
      that careful construction \emph{improves} performance rather than
      inflating it, contrary to the usual warning.
\item A pre-registered decision rule, stated before results were seen, for
      whether to escalate to representation learning, together with the
      measurement that answered it.
\item A quantified account of what limits the task, expressed as the label
      supply a future collection would need.
\end{enumerate}

\section{Data and units of analysis}
\label{sec:data}

\subsection{Source}

We use the original Dryad deposit rather than the merged convenience files in
common circulation. The merged version resamples every channel onto a shared
grid, which discards the blood volume pulse and inter-beat interval streams and
therefore removes the only route to heart rate variability. Channels are read
at their native rates: acceleration at \SI{32}{\hertz}, blood volume pulse at
\SI{64}{\hertz}, electrodermal activity and temperature at \SI{4}{\hertz}, and
heart rate at \SI{1}{\hertz}. These rates are read from the file headers.
Published descriptions giving \SI{72}{\hertz} and \SI{10}{\hertz} for pulse and
temperature disagree with the files themselves, and we follow the files.

The archive contains 609 recording sessions totalling \SI{1251.7}{\hour} across
250 person-days. The median session is \SI{1.17}{\hour} rather than a full
shift, so a nurse's week is fragmented across roughly 40 short recordings.

\subsection{Clock alignment}

Sensor timestamps are UTC; survey timestamps are bare local wall-clock with no
zone recorded. Choosing wrongly places nearly every reported episode outside
its own recording and raises no error. We resolved it empirically by counting,
under each hypothesis, how many rated episodes start inside a session belonging
to the same subject. Treating survey times as UTC places 50 of 245 (20.4\%);
treating them as naive \texttt{America/Chicago} places 212 (86.5\%). A fixed
five-hour offset reaches 198, so the daylight-saving-aware zone is required:
the archive spans April to December, and a fixed offset misaligns the winter
episodes by an hour. Thirty-three rated episodes (13.5\%) fall outside every
session even under the winning hypothesis and are set aside.

\subsection{Window and label}

The unit of prediction is a \SI{120}{\second} window with a \SI{60}{\second}
hop. Two arguments fix this length. No labelled episode is shorter than
\SI{120}{\second}, and a \SI{300}{\second} window would discard 15 episodes
outright. Independently, ultra-short heart rate variability studies find
\SI{120}{\second} sufficient for time-domain measures
\cite{munoz2015}, with root mean square of successive differences validated
from recordings as short as \SI{15}{\second} \cite{orini2023}; the classical
five-minute figure is a recording convention rather than a requirement.

A window is \textbf{positive} if at least half of it falls inside a
high-severity episode. A window is \textbf{negative} only if it survives the
construction in Section \ref{sec:negatives}. Everything else is
\textbf{excluded}, not set to zero.

\subsection{The structural problem}

The survey contains 358 rows, of which 245 carry a severity rating, split
46 low, 20 medium, 179 high. \textbf{None of them means ``not stressed.''}
Every row is an episode a nurse actively flagged, and the rating is its
severity. The majority labelled class is therefore the most severe one, and
nine of fifteen subjects have no medium-severity episodes at all.

Two consequences follow. Three-class modelling is unworkable, so we collapse to
high-severity against a constructed baseline. And the comparison group does not
exist in the data, which makes this a positive-unlabelled problem
\cite{elkan2008} rather than ordinary binary classification.

\subsection{Exclusions and cohort}

Exclusions are applied in a declared order, because they do not commute:
measuring session length before removing non-wear time retains 517 sessions,
and measuring it after retains 504. Thirteen sessions turn on rule order alone.

We remove sessions under five minutes, non-wear windows, sessions with a dead
electrodermal sensor, unrated episodes from both classes, three duplicate
survey rows, and low- and medium-severity episodes from the binary task. We
merge twelve overlapping episode pairs to their outer span.

We further remove five subjects. One has no eligible comparison windows at all.
Four have median skin conductance between \SI{0.07}{} and
\SI{0.11}{\micro\siemens}, at or below the sensor floor. A low median alone
does not disqualify a subject, since a small but varying signal remains
informative after normalisation. The deciding statistic is relative variation,
the interquartile range divided by the median: these four score 0.44 to 0.72,
while the remaining eleven score 1.04 to 5.14, with no overlap. Their signal is
flat, not merely small.

The retained cohort is 10 subjects, 24{,}334 windows, and 138 high-severity
episodes with sensor coverage. Three counts describe this dataset and only two
are denominators: the window count is a compute statistic, the honest
denominator for detecting an episode is 138, and for generalising to a new
person it is 10.

\section{Do the labels mark anything physiological?}
\label{sec:labels}

Before building a classifier we ask whether there is anything to classify. We
align every high-severity episode at its reported onset, normalise each channel
against the subject's own preceding hour, and compare the first ten minutes
after onset against the preceding thirty.

\begin{table}[t]
\caption{Event-triggered response at reported onset, 140 episodes with usable
coverage. Robust statistics only; see the note below on why.}
\label{tab:onset}
\centering
\begin{tabular}{lrrrr}
\toprule
Channel & Median $\Delta z$ & Episodes positive & Sign test & Median raw shift \\
\midrule
Electrodermal activity & $+0.568$ & 90/140 (64.3\%) & $p = 0.0009$ &
  $+0.64\,\si{\micro\siemens}$ \\
Heart rate & $-0.067$ & 68/140 (48.6\%) & $p = 0.80$ & $+0.47$ bpm \\
\bottomrule
\end{tabular}
\end{table}

Table \ref{tab:onset} gives the result. Skin conductance rises at onset in
64.3\% of episodes, a shift that survives both a sign test and a Wilcoxon
signed-rank test ($p = 3.3 \times 10^{-6}$), and the raw median shift of
\SI{0.64}{\micro\siemens} sits on a baseline of \SI{0.68}{\micro\siemens},
roughly a doubling. Heart rate rises in 48.6\% of episodes, which is
indistinguishable from chance.

\paragraph{Why we report medians.} The pooled mean tells a different and false
story. It reports an electrodermal shift of $+9.75$ normalised units, of which
63\% is contributed by five episodes out of 140. The cause is a numerical
failure in the normalisation rather than a physiological one: the trailing
interquartile range used as a denominator approaches zero during flat stretches,
and the largest single value reached $+1600$ where a plausible score is near
$\pm 3$. We bound the denominator at its empirical first percentile and clip
the output, and we fix the use of medians and sign tests in advance, because
the same data yields opposite headlines otherwise.

This result reorders the rest of the study. The problem is well posed, and it
is well posed in one channel. Any subsequent result driven by heart rate
deserves suspicion until explained.

\section{Constructing a comparison group}
\label{sec:negatives}

\subsection{Eligibility}

A candidate negative window must satisfy four conditions: it lies in a session
of at least thirty minutes; it falls on a subject-day containing at least one
reported episode; it lies at least thirty minutes from the boundary of any
reported episode, including unrated ones; and it is not itself positive.

The second condition carries most of the argument. On a day when a nurse filed
nothing, her silence is uninformative: she may have been calm, or too busy to
open the application. On a day when she filed at least one report, she was
demonstrably engaging with the survey, so silence during the remainder of that
day is evidence rather than absence of evidence. This condition is also the
most expensive, reducing 24{,}334 windows to 16{,}214. All four together leave
7{,}955 candidate windows, or \SI{265.2}{\hour}.

\subsection{Matching rather than filtering}

Nurses move more when busy, and busy periods are stressful, so movement
correlates with stress without being a stress signal. Selecting quiet periods
as negatives would produce two classes differing in both stress and movement,
and a classifier would separate them by detecting movement.

We therefore match rather than filter. We bin the positive class into deciles
of mean acceleration magnitude, and for each bin draw negatives from the same
bin, \emph{within the same subject}. Global matching would draw negatives
disproportionately from high-coverage subjects, making subject identity a
partial label proxy along exactly the axis the evaluation holds out.

\subsection{The negative ratio}

We fix the negative-to-positive ratio uniformly rather than letting each
subject supply its maximum. The reason is that class balance determines where
the decision threshold belongs: a fixed model with recall 0.80 and specificity
0.80 attains $F_1 = 0.727$ in a fold that is one-third positive and
$F_1 = 0.195$ in a fold that is one-thirtieth positive. Per-fold metrics would
not be comparable.

Exact uniformity is unattainable. The binding constraint is a single cell, one
subject's highest-activity bin containing 13 positives and 7 candidates, so
matching every subject exactly would require a ratio below one. We report a
sensitivity analysis across four arms in Table \ref{tab:sens} and fix
$\text{ratio} = 1.31$, which halves the residual spread from
$1.64\times$ to $1.28\times$.

Negatives are drawn under three random seeds. The seed-to-seed spread is the
noise floor: any later comparison smaller than it is not a comparison.

\subsection{Falsification}

We then build models we want to fail. Each is given only one suspicious signal,
using the same learner as the baseline so that a null result is attributable to
absence of signal rather than to a weaker model.

\begin{table}[t]
\caption{Falsification probes and ratio sensitivity, mean of three seeds.
All probes pass in all arms. Area under the ROC curve is prevalence-invariant,
so arms with different denominators remain comparable.}
\label{tab:sens}
\centering
\begin{tabular}{lrrrrrr}
\toprule
Arm & Subjects & Balance spread & ACC only & Time only & EDA only & Full \\
\midrule
All 10, ratio 1.31   & 10 & $1.28\times$ & 0.513 & 0.507 & 0.768 & 0.749 \\
All 10, ratio 2.14   & 10 & $1.64\times$ & 0.511 & 0.510 & 0.770 & 0.736 \\
Drop 8B, ratio 1.31  &  9 & $1.07\times$ & 0.521 & 0.533 & 0.762 & 0.739 \\
Drop 8B, ratio 2.25  &  9 & $1.30\times$ & 0.517 & 0.525 & 0.765 & 0.726 \\
\bottomrule
\end{tabular}
\end{table}

A model given only acceleration reaches AUC 0.513, and one given only the hour
of day reaches 0.507, against a chance value of 0.500. Both hold in all four
arms of Table \ref{tab:sens}. The constructed comparison group therefore does
not leak movement or shift schedule.

Subject identity is recoverable from the same features at 0.768 accuracy
against a chance value of 0.100. We record this as a measurement rather than a
failure, since physiology genuinely differs between people; it is the reference
against which any learned representation must later be compared.

\section{Baseline}
\label{sec:baseline}

Features are extracted at native rates and then aggregated to the window.
Detecting skin conductance responses on a \SI{1}{\hertz} downsample loses
roughly 40\% of them, so the common \SI{1}{\hertz} grid is used only for the
label join and coarse features. The feature set comprises tonic level and
slope, phasic variability, response count and mean amplitude for electrodermal
activity; mean, variability, slope, maximum and deviation from a trailing
median for heart rate; magnitude statistics and a movement fraction for
acceleration; slope and deviation from a trailing median for temperature, never
its absolute value, which has an intraclass correlation of 0.52 with subject
identity and acts as a fingerprint; and hour of day with position in session.

Heart rate variability is excluded from the baseline. Although the inter-beat
interval file is present in every session, only 5.9\% of \SI{120}{\second}
windows contain a run of clean consecutive beats long enough to compute it, and
the archive holds only \SI{39.5}{\hour} of usable beat data out of
\SI{1251.7}{\hour}. Reported session-level availability of 48\% is eight times
more flattering than the window-level figure that governs a per-window feature.

Continuous channels are normalised against a trailing sixty-minute robust
baseline within session, which uses only past data and is therefore both
leakage-free and reproducible at deployment.

The model is gradient boosting with balanced class weights, evaluated
leave-one-subject-out across three negative-sampling seeds.

\begin{table}[t]
\caption{Baseline performance and operating points. The alarm rate is
expressed against true worn time using the observed window prevalence of
6.49\%, not against the sampled negatives, which are enriched roughly
sevenfold in positives.}
\label{tab:baseline}
\centering
\begin{tabular}{lrrrr}
\toprule
Target alarms per worn hour & False positive rate & Window recall &
Episode recall & Seed range \\
\midrule
0.5 & 0.018 & 0.087 & 0.029 & 0.000 \\
1.0 & 0.036 & 0.155 & \textbf{0.106} & 0.072 \\
2.0 & 0.071 & 0.292 & 0.280 & 0.058 \\
\bottomrule
\end{tabular}
\end{table}

The model attains AUC $0.7654$ with a seed range of $0.0095$, and average
precision $0.685$ against a sampled prevalence of $0.436$. Table
\ref{tab:baseline} translates this into operating points.

\paragraph{The gap is the finding.} An AUC of 0.765 means the model reliably
ranks stressed windows above unstressed ones. At one false alarm per worn hour,
it recovers 10.6\% of episodes. AUC integrates over false positive rates
including regions no deployment would tolerate; once constrained to the usable
corner, little remains. Two of ten subjects detect zero episodes at any usable
threshold, and one is below chance at AUC 0.388.

We report per fold and never average. Fold AUC ranges from 0.388 to 0.885.

\section{Does a different treatment of the labels help?}
\label{sec:pu}

We pre-registered a decision rule before computing the following, and record it
here in the form it was written:

\begin{quote}
Proceed to representation learning only if a different label treatment moves
episode recall at one alarm per worn hour by more than the baseline seed range
of $0.0724$.
\end{quote}

The threshold is the baseline's own measured seed range rather than a chosen
constant. We compare three treatments on identical folds, windows, and
features: the curated negatives of Section \ref{sec:negatives}; a naive
comparator treating all unlabelled time as negative, with no same-day rule,
guard band, or matching; and non-negative positive-unlabelled risk estimation
\cite{kiryo2017} fitted subject-stratified, with the class prior bracketed
rather than estimated, since it is not identifiable from positive-unlabelled
data.

\begin{table}[t]
\caption{Label treatments on identical folds. Both reported metrics are
prevalence-invariant, so the naive arm scored over 24{,}334 windows and the
curated arm over 3{,}622 remain comparable.}
\label{tab:pu}
\centering
\begin{tabular}{llrr}
\toprule
Treatment & Class prior & AUC & Episode recall at 1 alarm/hour \\
\midrule
Curated negatives (baseline) & --- & \textbf{0.7654} & \textbf{0.1063} \\
Naive, all unlabelled negative & --- & 0.6925 & 0.0652 \\
Non-negative PU & 0.05 & 0.6616 & 0.0435 \\
Non-negative PU & 0.10 & 0.6933 & 0.0652 \\
Non-negative PU & 0.20 & 0.7132 & 0.0507 \\
Non-negative PU & 0.30 & 0.7167 & 0.0580 \\
\bottomrule
\end{tabular}
\end{table}

Table \ref{tab:pu} gives the result. The threshold to clear was $0.1787$. The
best alternative reached $0.0652$, a change of $-0.0411$. \textbf{The rule was
not met, and every alternative is worse than the baseline rather than merely
insufficiently better.}

\paragraph{Curation improved rather than inflated.} This comparison is normally
run to detect inflation: naive negatives include off-shift and sleeping time,
the contrast becomes trivial, and the score rises. Here the curated
construction beats naive by 0.073 AUC and 0.041 episode recall. The likely
mechanism is the guard band, since naive negatives include windows adjacent to
episodes, which are physiologically similar to positives and act as label noise
during training. The same-day rule and guard band are therefore not only
defensibility measures; they measurably improve the model.

\paragraph{The prior does not rescue it.} AUC rises monotonically with the
class prior, from 0.662 to 0.717, which is the shape one could present as
promising. Episode recall at the deployable operating point does not follow it,
remaining between 0.044 and 0.065 throughout. We report the prior-indexed
results in full rather than the trend alone.

\section{Discussion}

\subsection{What the eliminations leave}

The labels mark a real physiological state, in one channel. The comparison
group does not leak movement or schedule. Standard features and a standard
learner recover one episode in ten at a deployable alarm rate. Three distinct
treatments of the labels do not improve on this, and two make it worse.

What remains is the data. There are 138 episodes from 10 subjects; the labels
are retrospective self-reports of intervals, treated here as uniform states
although a thirty-minute report is unlikely to describe thirty uniform minutes;
the negative class is constructed rather than observed; and only one of four
channels responds at onset.

\subsection{On the reported accuracy in prior work}

Mathur et al.\ \cite{mathur2024} report weighted $F_1 = 0.99$ on this dataset.
We do not reproduce this and we think it is difficult to reconcile with
subject-grouped evaluation. Adjacent windows in these recordings share raw
samples and subject-specific baselines; a split that does not respect subject
boundaries places near-duplicates on both sides. We are not in a position to
diagnose another group's protocol from a published summary, and we note the
discrepancy rather than attributing a cause. We report our own protocol in full
so that the comparison can be made.

\subsection{What a future collection would need}

The most actionable consequence concerns label supply rather than modelling.
A collection targeting this task should record explicit non-stress intervals,
since their absence forces every downstream number to be conditional on a
construction; should capture onset with finer resolution than a retrospective
interval; and should prioritise electrodermal signal quality, since four of
fifteen subjects here produced a flat trace on the only responsive channel.
Heart rate contributed nothing measurable at onset in this cohort, and
inter-beat intervals were usable in under 6\% of windows.

\section{Limitations}

The cohort is 10 of 15 female nurses at a single hospital during a pandemic,
and external validity is correspondingly narrow. Negatives are constructed, so
precision is not reported as a performance claim: when the model flags an
unlabelled window we cannot distinguish a false alarm from an unreported
episode. The class prior is not identifiable and is bracketed rather than
estimated. Our non-negative PU implementation uses a fixed training budget
without an inner validation split, so its negative result is weaker evidence
than the naive comparator, which shares the baseline's learner and features and
differs only in label treatment. Thirty-three rated episodes have no sensor
coverage and are absent from every metric. Two normalisation constants, the
denominator floor and the output clip, are choices we record and vary rather
than defend. Ten folds is underpowered for the paired comparisons we report,
and we state spreads descriptively rather than dividing them by the square root
of the fold count.

\section{Conclusion}

We set out to determine whether the ceiling on wearable stress detection in
this dataset lies in the method or in the data, and answered by elimination.
The labels are real, the comparison group is sound, the baseline is honest, and
three ways of using the labels differently do not move it. The binding
constraint is the label supply. We report this as a result rather than a
failure, because the quantity a future study most needs from this one is not a
score but a measurement of what was missing.

\begin{ack}
We thank the authors of the source dataset for making it publicly available.
\end{ack}

\begin{thebibliography}{9}

\bibitem{elkan2008}
C.~Elkan and K.~Noto.
\newblock Learning classifiers from only positive and unlabeled data.
\newblock In \emph{Proceedings of the 14th ACM SIGKDD International Conference
  on Knowledge Discovery and Data Mining}, 2008.
\newblock \doi{10.1145/1401890.1401920}.

\bibitem{hosseini2022}
S.~Hosseini, R.~Gottumukkala, S.~Katragadda, R.~T. Bhupatiraju, Z.~Ashkar,
  C.~W. Borst, and K.~Cochran.
\newblock A multimodal sensor dataset for continuous stress detection of nurses
  in a hospital.
\newblock \emph{Scientific Data}, 9:255, 2022.
\newblock \doi{10.1038/s41597-022-01361-y}.

\bibitem{kiryo2017}
R.~Kiryo, G.~Niu, M.~C. du~Plessis, and M.~Sugiyama.
\newblock Positive-unlabeled learning with non-negative risk estimator.
\newblock In \emph{Advances in Neural Information Processing Systems 30}, 2017.
\newblock arXiv:1703.00593.

\bibitem{mathur2024}
A.~Mathur et~al.
\newblock Body sensor-based multimodal nurse stress detection using machine
  learning.
\newblock In \emph{2024 16th International Conference on Communication Systems
  and Networks (COMSNETS)}, 2024.

\bibitem{munoz2015}
M.~L. Munoz, A.~van~Roon, H.~Riese, C.~Thio, E.~Oostenbroek, I.~Westrik,
  E.~J.~C. de~Geus, R.~Gansevoort, J.~Lefrandt, I.~M. Nolte, and H.~Snieder.
\newblock Validity of (ultra-)short recordings for heart rate variability
  measurements.
\newblock \emph{PLOS ONE}, 10(9):e0138921, 2015.
\newblock \doi{10.1371/journal.pone.0138921}.

\bibitem{orini2023}
M.~Orini et~al.
\newblock Long-term association of ultra-short heart rate variability with
  cardiovascular events.
\newblock \emph{Scientific Reports}, 13, 2023.

\end{thebibliography}

\end{document}
